\begin{longtable}{ | L{0.15\textwidth} | L{0.45\textwidth} | L{0.25\textwidth} | }
    \hline
    \textbf{Notation} & \textbf{Explanation} & \textbf{Purpose} \\
    \hline
    \endfirsthead
    
    \hline
    \textbf{Notation} & \textbf{Explanation} & \textbf{Purpose} \\
    \hline
    \endhead

    \hline
    \endfoot

    \hline
    \endlastfoot

    % Entries start here
    $x \in A$ 
    & Element $x$ belongs to set $A$.
    & Basic set membership. \\
    \hline

    $A \subseteq B$
    & Set $A$ is a subset of set $B$.
    & Indicates inclusion. \\
    \hline

    $A \cap B$
    & Intersection of sets $A$ and $B$ (elements common to both).
    & Logical AND operation on sets. \\
    \hline

    $A \cup B$
    & Union of sets $A$ and $B$ (elements in either $A$ or $B$).
    & Logical OR operation on sets. \\
    \hline

    $\emptyset$
    & The empty set containing no elements.
    & Identity element for union. \\
    \hline

    $2^A$ or $\mathcal{P}(A)$
    & Power set of $A$ (set of all subsets of $A$).
    & Defining sigma-algebras. \\
    \hline
    
    $\forall$ 
    & Universal Quantifier (For all) 
    & Defining properties that must hold for every element (e.g., continuity). \\ \hline

    $\exists$ 
    & Existential Quantifier (There exists) 
    & Asserting the existence of a specific object (e.g., existence of a limit). \\ \hline

    $f^{-1}(B)$ 
    & Inverse Image 
    & The core tool for defining Random Variables (Mapping values back to events). \\ \hline

\end{longtable}
